\documentclass[a4paper]{article}
\usepackage[utf8]{inputenc}
\usepackage{amsmath, amssymb, amsfonts}
\usepackage{enumitem}
\pagenumbering{gobble}

\begin{document}

\begin{center}
  \textbf{\underline{Quiz 1 Teori Approksimasi Gasal 2025/2026}}\\[1em]
  \begin{tabular}{lcl}
    Mata Kuliah/SKS & : & Teori Approksimasi / 3 SKS \\
    Hari, Tanggal   & : & Rabu, 10 September 2025    \\
    Waktu           & : & 09.00-10.40 WIB            \\
    Sifat           & : & Terbuka
  \end{tabular}
\end{center}

\vspace{1.5em}

\begin{enumerate}
  \item Show that
        \[
          \left| e^x - \sum_{n=0}^8 \frac{x^n}{n!} \right| \le 2.8 \times 10^{-10}, \forall x \in [-1, 0].
        \]

  \item \begin{enumerate}[label=(\alph*), before=\vspace{-\dimexpr\baselineskip+\topsep\relax}]
          \item Find a Polynomial $P(x) = \sum\limits_{n=0}^N a_n x^n$ such that
                \[
                  \left| \sin x - \sum_{n=0}^N a_n x^n \right| \le 0.0001, \forall x \in \left[0, \frac{\pi}{2}\right].
                \]
          \item Use the result to find an approximative value of
                \[
                  \int_0^1 \sin x^4 dx,
                \]
                with an error of at most 0.0001.
        \end{enumerate}

  \item Dapatkan radius konvergensi untuk deret pangkat
        \[
          \sum_{n=0}^\infty \frac{n+2}{n+1} x^{n+1}.
        \]

  \item Buktikan bahwa
        \[
          \sum_{n=k+1}^\infty \frac{1}{(2n-1)^2} \le \frac{1}{4k+2} + \frac{1}{(2k+1)^2}, \forall k \in \mathbb{N}.
        \]
\end{enumerate}

\newpage

\begin{center}
  \textbf{\underline{Solusi Quiz 1 Teori Approksimasi Gasal 2025/2026}}
\end{center}

\begin{enumerate}
  \item Pernyataan di atas \textbf{tidak benar} (salah). Untuk menyangkal suatu pernyataan yang berkuantor universal ($\forall x \in [-1, 0]$), cukup ditunjukkan satu contoh penyangkal (\textit{counterexample}) pada interval tersebut di mana ketaksamaan tidak terpenuhi.

        Pilihlah titik ujung $x = -1 \in [-1, 0]$.

        Misalkan $P_8(x) = \sum_{n=0}^8 \frac{x^n}{n!}$ adalah polinomial Taylor (Maclaurin) berderajat $8$ dari fungsi $f(x) = e^x$. Nilai aproksimasi pada $x = -1$ adalah:
        \begin{align*}
          P_8(-1) & = \sum_{n=0}^8 \frac{(-1)^n}{n!}                                                                                      \\
                  & = 1 - 1 + \frac{1}{2!} - \frac{1}{3!} + \frac{1}{4!} - \frac{1}{5!} + \frac{1}{6!} - \frac{1}{7!} + \frac{1}{8!}      \\
                  & = 1 - 1 + \frac{1}{2} - \frac{1}{6} + \frac{1}{24} - \frac{1}{120} + \frac{1}{720} - \frac{1}{5040} + \frac{1}{40320} \\
                  & = \frac{14833}{40320} \approx 0{,}3678819444.
        \end{align*}
        Sementara itu, nilai eksak dari $e^{-1} = \frac{1}{e}$ hingga 10 desimal adalah:
        \[
          e^{-1} \approx 0,3678794412.
        \]
        Dengan demikian, galat aproksimasi aktual pada titik $x = -1$ bernilai:
        \[
          \left| e^{-1} - \sum_{n=0}^8 \frac{(-1)^n}{n!} \right| = |0,3678794412 - 0,3678819444| \approx 2,50327 \times 10^{-6}.
        \]
        Karena
        \[
          2,50327 \times 10^{-6} > 2,8 \times 10^{-10},
        \]
        maka ketaksamaan tersebut \textbf{gagal} pada $x = -1 \in [-1, 0]$.

  \item \begin{enumerate}[label=(\alph*)]
          \item Deret Maclaurin untuk fungsi $f(x) = \sin x$ diberikan oleh
                \[
                  \sin x = \sum_{k=0}^\infty \frac{(-1)^k}{(2k+1)!} x^{2k+1} = x - \frac{x^3}{3!} + \frac{x^5}{5!} - \frac{x^7}{7!} + \frac{x^9}{9!} - \dots
                \]
                Kita pilih polinomial aproksimasi berupa polinomial Taylor berderajat $N = 2m+1$:
                \[
                  P(x) = P_{2m+1}(x) = \sum_{k=0}^m \frac{(-1)^k}{(2k+1)!} x^{2k+1}.
                \]
                Karena turunan ke-$(2m+2)$ dari $\sin x$ di titik $x = 0$ bernilai nol, polinomial Taylor berderajat ganjil $2m+1$ identik dengan polinomial Taylor berderajat genap $2m+2$, yaitu $P_{2m+1}(x) = P_{2m+2}(x)$.

                Berdasarkan Teorema Taylor bentuk Lagrange, untuk setiap $x \in [0, \pi/2]$, terdapat $\xi \in (0, x)$ sedemikian sehingga sisa aproksimasi memenuhi
                \[
                  |R_{2m+2}(x)| = \left| \frac{f^{(2m+3)}(\xi)}{(2m+3)!} x^{2m+3} \right|.
                \]
                Setiap turunan dari fungsi sinus bernilai $\pm \cos x$ atau $\pm \sin x$, sehingga $|f^{(2m+3)}(\xi)| \le 1$. Karena untuk $x \in [0, \pi/2]$ berlaku $|x| \le \frac{\pi}{2}$, maka batas galat maksimum adalah
                \[
                  |R_{2m+2}(x)| \le \frac{(\pi/2)^{2m+3}}{(2m+3)!}.
                \]
                Kita uji nilai derajat polinomial:
                \begin{itemize}
                  \item Untuk $m = 3$ (derajat $N = 7$):
                        \[
                          \frac{(\pi/2)^9}{9!} \approx \frac{(1{,}570796)^9}{362.880} \approx \frac{58{,}2189}{362.880} \approx 1{,}604 \times 10^{-4} > 0{,}0001.
                        \]
                        (Belum mencukupi ketelitian ${,}0001$).
                  \item Untuk $m = 4$ (derajat $N = 9$):
                        \[
                          \frac{(\pi/2)^{11}}{11!} \approx \frac{(1{,}570796)^{11}}{39.916.800} \approx \frac{143{,}646}{39.916.800} \approx 3{,}599 \times 10^{-6} \le 0{,}0001.
                        \]
                        (Bahkan jika ditinjau dari rumus sisa berderajat $9$, batas galat $\frac{(\pi/2)^{10}}{10!} \approx 2{,}54 \times 10^{-5} < 0{,}0001$).
                \end{itemize}
                Dengan demikian, derajat terkecil yang menjamin galat $\le 0{,}0001$ adalah $N = 9$. Polinomial $P(x)$ yang dicari adalah
                \begin{align*}
                  P(x) & = x - \frac{x^3}{3!} + \frac{x^5}{5!} - \frac{x^7}{7!} + \frac{x^9}{9!}        \\
                       & = x - \frac{x^3}{6} + \frac{x^5}{120} - \frac{x^7}{5040} + \frac{x^9}{362880}.
                \end{align*}

          \item Dari hasil bagian (a), untuk setiap $u \in [0, \pi/2]$, berlaku
                \[
                  |\sin u - P(u)| \le 0{,}0001.
                \]
                Untuk setiap $x \in [0, 1]$, perhatikan bahwa $x^4 \in [0, 1]$. Karena $[0, 1] \subset [0, \pi/2]$ (mengingat $\pi/2 \approx 1{,}5708 > 1$), maka kita dapat mensubstitusikan $u = x^4$ ke dalam ketaksamaan tersebut:
                \[
                  |\sin(x^4) - P(x^4)| \le 0{,}0001, \quad \forall x \in [0, 1].
                \]
                Dengan mengintegralkan kedua ruas pada interval $[0, 1]$, kita peroleh
                \begin{align*}
                  \left| \int_0^1 \sin(x^4) \, dx - \int_0^1 P(x^4) \, dx \right|
                   & \le \int_0^1 \left| \sin(x^4) - P(x^4) \right| dx \\
                   & \le \int_0^1 0{,}0001 \, dx                       \\
                   & = {,}0001 \times (1 - 0) = {,}0001.
                \end{align*}
                Oleh karena itu, nilai aproksimasi integral yang dicari adalah
                \begin{align*}
                  \int_0^1 P(x^4) \, dx
                   & = \int_0^1 \left( x^4 - \frac{x^{12}}{3!} + \frac{x^{20}}{5!} - \frac{x^{28}}{7!} + \frac{x^{36}}{9!} \right) dx               \\
                   & = \left[ \frac{x^5}{5} - \frac{x^{13}}{78} + \frac{x^{21}}{2520} - \frac{x^{29}}{146160} + \frac{x^{37}}{13426560} \right]_0^1 \\
                   & = \frac{1}{5} - \frac{1}{78} + \frac{1}{2520} - \frac{1}{146160} + \frac{1}{13426560}.
                \end{align*}
                Nilai dari masing-masing suku pecahan tersebut adalah:
                \begin{align*}
                  \frac{1}{5}        & = {,}2000000000         \\
                  -\frac{1}{78}      & \approx -0{,}0128205128 \\
                  \frac{1}{2520}     & \approx +0{,}0003968254 \\
                  -\frac{1}{146160}  & \approx -0{,}0000068418 \\
                  \frac{1}{13426560} & \approx +0{,}0000000745
                \end{align*}
                Menjumlahkan seluruh suku menghasilkan nilai aproksimasi
                \[
                  \int_0^1 \sin(x^4) \, dx \approx 0{,}1875695.
                \]
                \textit{Catatan:} Deret ini merupakan deret bolak-balik (\textit{alternating series}) yang suku-sukunya monoton turun menuju nol. Suku ke-4 adalah $\frac{1}{146160} \approx 6{,}84 \times 10^{-6} < 0{,}0001$. Oleh karena itu, bahkan dengan menggunakan 3 suku pertama saja:
                \[
                  \frac{1}{5} - \frac{1}{78} + \frac{1}{2520} = \frac{1229}{6552} \approx 0{,}187576,
                \]
                galat aproksimasi sudah dijamin kurang dari $\frac{1}{146160} < 0{,}0001$.
        \end{enumerate}

  \item Misalkan suku ke-$n$ dari deret tersebut adalah
        \[
          u_n = \frac{n+2}{n+1} x^{n+1}.
        \]
        Berdasarkan Uji Rasio (\textit{Ratio Test}) untuk konvergensi mutlak, kita evaluasi limit
        \[
          L = \lim_{n \to \infty} \left| \frac{u_{n+1}}{u_n} \right|.
        \]
        Suku ke-$(n+1)$ adalah
        \[
          u_{n+1} = \frac{(n+1)+2}{(n+1)+1} x^{(n+1)+1} = \frac{n+3}{n+2} x^{n+2}.
        \]
        Rasio kedua suku tersebut adalah
        \begin{align*}
          \left| \frac{u_{n+1}}{u_n} \right|
           & = \left| \frac{\frac{n+3}{n+2} x^{n+2}}{\frac{n+2}{n+1} x^{n+1}} \right| \\
           & = \frac{(n+3)(n+1)}{(n+2)^2} \cdot |x|                                   \\
           & = \frac{n^2 + 4n + 3}{n^2 + 4n + 4} \cdot |x|.
        \end{align*}
        Dengan mengambil limit untuk $n \to \infty$:
        \[
          L = \lim_{n \to \infty} \left( \frac{n^2 + 4n + 3}{n^2 + 4n + 4} \right) |x| = \lim_{n \to \infty} \left( \frac{1 + \frac{4}{n} + \frac{3}{n^2}}{1 + \frac{4}{n} + \frac{4}{n^2}} \right) |x| = 1 \cdot |x| = |x|.
        \]
        Berdasarkan Uji Rasio:
        \begin{itemize}
          \item Deret konvergen mutlak jika $L < 1 \iff |x| < 1$.
          \item Deret divergen jika $L > 1 \iff |x| > 1$.
        \end{itemize}
        Karena deret berpusat di $x_0 = 0$ dan konvergen mutlak untuk $|x| < 1$, maka sesuai dengan definisi pertidaksamaan $|x - x_0| < R$, radius konvergensinya adalah $R = 1$.


  \item  Perhatikan suku pertama deret pada ruas kiri, yaitu pada saat indeks $n = k+1$:
        \[
          \text{Untuk } n = k+1: \quad \frac{1}{(2(k+1)-1)^2} = \frac{1}{(2k+2-1)^2} = \frac{1}{(2k+1)^2}.
        \]
        Dengan memisahkan suku pertama tersebut dari penjumlahan, ruas kiri dapat ditulis sebagai
        \[
          \sum_{n=k+1}^\infty \frac{1}{(2n-1)^2} = \frac{1}{(2k+1)^2} + \sum_{n=k+2}^\infty \frac{1}{(2n-1)^2}.
        \]
        Oleh sebab itu, pertidaksamaan pada soal ekuivalen dengan membuktikan bahwa
        \[
          \sum_{n=k+2}^\infty \frac{1}{(2n-1)^2} \le \frac{1}{4k+2}.
        \]
        Untuk membuktikan hal ini, kita gunakan Teorema Uji Integral untuk Estimasi Sisa Deret.
        Definisikan fungsi
        \[
          f(x) = \frac{1}{(2x-1)^2}, \quad \text{untuk } x \ge 1.
        \]
        Fungsi $f(x)$ kontinu, positif, dan monoton turun pada interval $[1, \infty)$ karena turunan pertamanya adalah
        \[
          f'(x) = -\frac{4}{(2x-1)^3} < 0, \quad \forall x > \frac{1}{2}.
        \]
        Karena $f(x)$ monoton turun, untuk setiap interval $[n-1, n]$ dengan $n \ge k+2$, berlaku
        \[
          f(n) \le f(x), \quad \forall x \in [n-1, n].
        \]
        Dengan mengintegralkan kedua ruas pada interval $[n-1, n]$ yang panjangnya $\Delta x = 1$:
        \[
          \frac{1}{(2n-1)^2} = f(n) = \int_{n-1}^n f(n) \, dx \le \int_{n-1}^n f(x) \, dx = \int_{n-1}^n \frac{1}{(2x-1)^2} \, dx.
        \]
        Jumlahkan ketaksamaan ini dari $n = k+2$ sampai $\infty$:
        \[
          \sum_{n=k+2}^\infty \frac{1}{(2n-1)^2} \le \sum_{n=k+2}^\infty \int_{n-1}^n \frac{1}{(2x-1)^2} \, dx = \int_{k+1}^\infty \frac{1}{(2x-1)^2} \, dx.
        \]
        Sekarang kita hitung integral tak wajar di ruas kanan:
        \begin{align*}
          \int_{k+1}^\infty \frac{1}{(2x-1)^2} \, dx
           & = \lim_{b \to \infty} \int_{k+1}^b (2x-1)^{-2} \, dx = \lim_{b \to \infty} \left[ -\frac{1}{2(2x-1)} \right]_{k+1}^b                                                                                     \\
           & = \lim_{b \to \infty} \left( -\frac{1}{2(2b-1)} \right) - \left( -\frac{1}{2(2(k+1)-1)} \right)                                                                                                          \\
           & = 0 + \frac{1}{2(2k+2-1)}                                                                                            = \frac{1}{2(2k+1)}                                               = \frac{1}{4k+2}.
        \end{align*}
        Dengan demikian, terbukti bahwa
        \[
          \sum_{n=k+2}^\infty \frac{1}{(2n-1)^2} \le \frac{1}{4k+2}.
        \]
        Terakhir, dengan menambahkan suku $\frac{1}{(2k+1)^2}$ pada kedua ruas, kita peroleh
        \[
          \sum_{n=k+1}^\infty \frac{1}{(2n-1)^2} = \frac{1}{(2k+1)^2} + \sum_{n=k+2}^\infty \frac{1}{(2n-1)^2} \le \frac{1}{4k+2} + \frac{1}{(2k+1)^2}.
        \]
        Pertidaksamaan ini berlaku untuk setiap $k \in \mathbb{N}$.
        \hfill $\blacksquare$
\end{enumerate}

\end{document}
